\documentclass[12pt,a4paper]{article}
\usepackage[utf8]{inputenc}
\usepackage[T1]{fontenc}
\usepackage{amsmath,amssymb,amsfonts}
\usepackage{booktabs}
\usepackage{hyperref}
\usepackage[margin=2.5cm]{geometry}
\usepackage[numbers]{natbib}

\newcommand{\SBD}{S_{\mathrm{BD}}}
\newcommand{\Lor}{\mathrm{Lor4D}}

\title{Curvature Encoding in Finite Causal Sets:\\
  Purely Causal Observables Across de~Sitter, Schwarzschild, and FLRW Backgrounds}

\author{Gang Zhang}
\date{\today}

\begin{document}
\maketitle

\begin{abstract}
We present numerical evidence that finite causal sets encode background
spacetime curvature through observables constructed solely from the partial
order. Across 23 experiments on de~Sitter, Schwarzschild, and FLRW backgrounds
at scales $N=16$--$1024$, we find three recurring features. First, the
occupancy ratio exhibits a sharp gate-like response as curvature increases.
Second, the antichain-based channel tracks the Hubble parameter more robustly
than the spectral channel across all tested backgrounds, including local
$H(t)$ variation in FLRW examples. Third, a calibrated two-step construction
provides a high-fidelity bridge from causal observables to scalar curvature at
the largest finite sizes considered, although the bridge remains sensitive to
finite-$N$ stiffness.

The evidence is consistent with the Benincasa--Dowker action carrying genuine
curvature information at finite $N$ and with antichain observables being the
dominant finite-size curvature carrier in the tested backgrounds. Our results
are numerical and finite-size in nature; analytic control of the second-order
bridge remains an open problem.
\end{abstract}

\noindent\textbf{Keywords:}
causal sets, curvature encoding, Benincasa--Dowker action, Einstein--Hilbert
action, de Sitter, Schwarzschild, FLRW, partial order observables

% ============================================================
\section{Introduction}
% ============================================================

A foundational question in causal set theory is whether the discrete causal
structure of a finite poset contains enough information to recover the
continuous geometry of the spacetime from which it was sprinkled.
The Benincasa--Dowker (BD) action~\cite{BenincasaDowker2010} provides a
partial answer: it reproduces the Einstein--Hilbert action in the continuum
limit for manifold-like causal sets.
But does this correspondence hold at \emph{finite} $N$, and does it extend
beyond the BD action to other causal observables?

This paper presents systematic evidence that finite causal sets encode
background curvature through purely causal observables that do not use linear
extension counts.

The main physical question is whether the Benincasa--Dowker action and related
causal observables retain curvature sensitivity at finite $N$, beyond the
continuum limit in which the BD action is usually discussed~\cite{BenincasaDowker2010,Surya2019}.

% ============================================================
\section{Experimental Setup}
% ============================================================

\subsection{Background geometries}

We generate Lorentzian causal sets by Poisson sprinkling into causal diamonds
in three background geometries:

\textbf{de~Sitter} ($d=4$): $a(t) = e^{Ht}$ with Hubble parameter
$H\in\{0, 0.01, 0.05, 0.1, 0.15, 0.2, 0.3, 0.5, 1.0, 2.0\}$.
The $H=0$ limit recovers flat Minkowski space.

\textbf{Schwarzschild} ($d=4$): weak-field approximation with compactness
parameter $\phi_0\in\{0, 0.01, 0.05, 0.1\}$.

\textbf{FLRW} ($d=4$): matter-dominated expansion $a(t)=(1+\kappa t)^{2/3}$
with $\kappa\in\{0, 0.3, 1.0, 3.0\}$.

Scales: $N\in\{16, 20, 28, 36, 48, 64, 96, 128, 256, 512, 1024\}$ with
15--40 realizations per $(N, \text{background})$ condition.

\subsection{Causal observables}

All observables are computed directly from the partial order, without any
reference to $\log H$:
\begin{itemize}
\item Occupancy ratio $R = C_0/\binom{N}{2}$ (fraction of comparable pairs)
\item Interval counts $C_k$ ($k=0,1,2,3$)
\item BD action $\SBD^{(4)} = N - C_0 + 9C_1 - 16C_2 + 8C_3$
\item Maximum antichain width $w$
\item Spectral gap of the causal matrix
\item Effective dimension $d_{\mathrm{eff}}$ (Myrheim--Meyer)
\end{itemize}

% ============================================================
\section{Three-Layer Curvature Encoding}
% ============================================================

\subsection{Layer 1: Sigmoid wall (curvature upper bound)}

The occupancy ratio $R$ exhibits a sigmoid response to background curvature:
as $H$ increases from 0, $R$ initially remains near its flat-space value, then
drops sharply beyond a critical curvature.
This sigmoid wall functions as a \emph{geometric gate}: posets sprinkled from
strongly curved backgrounds have $R$ values that fall outside the flat-space
range, making them identifiable as non-flat.

The wall position scales with $N$: at larger $N$, the wall sharpens and moves
to lower $H$, consistent with the expectation that larger causal sets resolve
curvature more finely.

\subsection{Layer 2: First-order bulk recovery of $H$}

Two independent observable channels track the Hubble parameter:

\textbf{Antichain channel}: The maximum antichain width $w/N$ decreases
monotonically with $H$ (curvature compresses the spatial cross-section).
DDT condition C2 is satisfied in 21/21 tested conditions.

\textbf{Spectral channel}: The spectral gap of the causal matrix increases
with $H$.
DDT condition C2 is satisfied in 6/18 conditions (weaker, more noise-sensitive).

The two channels converge to a single geometric degree of freedom at large $N$:
$|\rho(\text{antichain}, \text{spectral})|\to 0.86$, confirming that both are
measuring the same underlying curvature.

\textbf{Local $H(t)$ tracking}: In FLRW backgrounds where $H$ varies with time,
the antichain-based observable tracks the local expansion rate with 24/24
anti-monotone direction consistency.

\subsection{Layer 3: Second-order bridge to scalar curvature}

The Einstein--Hilbert bridge connects the BD action density to the scalar
curvature $\mathcal{R}$.
A direct fit $\SBD/N \sim \mathcal{R}$ has poor finite-$N$ convergence
($\alpha_{\mathrm{eff}}\to 2$ requires $N\sim 10^9$).

We identify a viable two-step construction:
\begin{enumerate}
\item Calibrate a first-order observable against $H$
  (e.g., $\hat{H} = f(w/N, N)$).
\item Square: $\hat{\mathcal{R}} = d(d-1)\hat{H}^2$.
\end{enumerate}

This yields $R^2 = 0.987$ (OLS) and $R^2 = 0.9996$ (rank-preserving fit)
at $N=1024$ in de~Sitter backgrounds.
The second-order bridge is the squared amplitude of the first-order deviation
from the flat baseline---not an independent observable channel.

% ============================================================
\section{Discussion}
% ============================================================

The present results point to a simple conclusion: finite causal sets retain
curvature information in several observables that can be constructed purely
from the partial order. The occupancy ratio responds sharply to background
curvature, the antichain channel provides the most stable finite-size tracking
of $H$, and the calibrated second-order construction produces a high-fidelity
bridge to scalar curvature at the largest sizes tested.

Among the observables considered here, the antichain channel appears to be the
dominant carrier of curvature information in the finite-$N$ regime. The spectral
channel is still useful, but it is noisier and less robust across the tested
backgrounds.

The Benincasa--Dowker action remains important as a benchmark because it is the
standard local causal-set observable with a direct continuum interpretation.
Our numerical results are consistent with the view that its finite-$N$ behavior
contains curvature-sensitive information, but they do not provide an analytic
derivation of that fact.

The main limitation of the present study is finite-size stiffness in the
second-order bridge. A more complete analytic understanding of why the bridge
converges as it does, and why antichain-based observables dominate the signal,
remains open.

% ============================================================
\section{A Note on Asymmetry Under Augmentation}
% ============================================================

A related but independent finding concerns the directionality of causal
structure growth.
In augmentation experiments (adding elements to an existing poset), the entropy
change is asymmetric: future-directed augmentation (adding to the top of the
Hasse diagram) produces larger entropy increase than past-directed augmentation
(adding to the bottom).

This structural asymmetry---future-like growth expands combinatorial freedom
more than past-like growth compresses it---suggests a possible structural time
arrow that does not rely on thermodynamic assumptions.
The effect is observed across multiple families and $N$ values, though with
substantial individual variation.

We present this as an independent auxiliary observation, not as part of the
curvature-encoding analysis.

% ============================================================
\section{Conclusion}
% ============================================================

We have shown numerically that finite causal sets encode background curvature
through observables built solely from the partial order. The occupancy ratio
shows a gate-like response to curvature, the antichain channel robustly tracks
the Hubble parameter, and a calibrated two-step construction yields a
high-fidelity bridge to scalar curvature at the finite sizes tested.

The principal limitation is the finite-$N$ stiffness of the second-order
bridge; analytic control of the large-$N$ approach remains open. These results
indicate that finite posets retain nontrivial curvature information and that
antichain-based observables are the strongest carriers of that signal in the
tested examples.

\begin{thebibliography}{9}
\bibitem{BenincasaDowker2010}
D.~M.~T.~Benincasa and F.~Dowker,
\emph{Scalar Curvature of a Causal Set},
\emph{Phys.\ Rev.\ Lett.}\ \textbf{104}, 181301 (2010),
arXiv:1001.2725.

\bibitem{Surya2019}
S.~Surya,
\emph{The causal set approach to quantum gravity},
\emph{Living Rev.\ Relativ.}\ \textbf{22}, 5 (2019),
arXiv:1903.11544.
\end{thebibliography}

\end{document}
